% page 420 of the facsimile (image p-419 of the scan)
% block 149.9 x 229.0 mm ; left margin 8.0 mm
\pgc{-12.700mm}{0.000mm}{
\l{\cel{3}412.}
\l{}
\l{\sou{orveto}.\cel{2}(\sou{zool.}) Mikra reptero sauria, ne-nociva, tipo}
\l{\cel{3}di la familio \textquotedbl{}anguidi\textquotedbl{}, qua vivas en Europa ed en}
\l{\cel{3}Azia orientala. - L.\cel{1}\sou{anguis fragilis}. - EF.}
\l{}
\l{\sou{orvietano}.Elektuario (varianto nenociva di teriako) quan}
\l{\cel{3}onu judikis, en la mezo di la 17esma yarcento, kom tre}
\l{\cel{3}apta e tre efikiva. - DEFIS.}
\l{}
\l{\sou{-os.}\cel{2}(\sou{gram.})\cel{2}Che la verbi : dezinenco di indikativo,}
\l{\cel{3}futuro. -}
\l{}
\l{\sou{osifrago}. (\sou{zool.)}\cel{2}Speco di raptucelo jornala, ek la fami-}
\l{\cel{3}lio \textquotedbl{}falkonidi\textquotedbl{}, qua vivas en la nordo-regioni di\cel{2}Eu-}
\l{\cel{3}ropa. - L.\cel{1}\sou{aquila ossifraga}. - EFI.}
\l{}
\l{\sou{oskular}. (\sou{trans.}) (\sou{geom.)}\cel{2}Efektigar kontakto supra-klasa}
\l{\cel{3}kun kurvo, surfaco, en un punto di altra kurvo, surfaco.}
\l{}
\l{\sou{osmio}. (\sou{kemio)}. Korpo simpla, metala, quan onu renkontras}
\l{\cel{3}en ula platin-erco. -\cel{1}\sou{Simbolo kemiala}\cel{1}: Os.- DEFIRS.}
\l{}
\l{\sou{osmoso.}\cel{2}I. (\sou{biol. e kem.})\cel{2}Fenomeno qua konsistas ye la}
\l{\cel{3}diplacesi qui eventas de ca a ta, ek du liquidi di qui}
\l{\cel{3}la koncentreso-gradi interdiferas, tra parieto mi-}
\l{\cel{3}permeebla. - DEFS.}
\l{}
\l{\sou{osto}.\cel{2}Parto harda e solida, ye qua konsistas la skeleto}
\l{\cel{3}che la animali vertebroza di la klasi supra. EFIS.}
\l{}
\l{\sou{ostensorio}. (\sou{liturgio katol.}) Suportilo ek arjento, arjen-}
\l{\cel{3}to orizita, oro, en qua onu expozas, a la adoro da la}
\l{\cel{3}religiani, la hostio sakrigita. - DEFI.}
\l{}
\l{\sou{ostentar}. (\sou{netrans., pri)}\cel{2}Montrar afektacoze to quo esas}
\l{\cel{3}apta flatar sua vanitato. - EFIS.}
\l{}
\l{\sou{osteoblasto.}\cel{1}(\sou{biol.)}\cel{2}Celulo yuna, restinta e retroveninta}
\l{\cel{3}a sua stando embrionala, quan onu vidas en la osto-}
\l{\cel{3}medulo e la strato ostifera di periosto. - DEFIRS.}
\l{}
\l{\sou{osteologio.}\cel{2}Parto di anatomio, qua traktas la osti. DEFIRS}
\l{}
\l{\sou{ostegomo}. (\sou{patol.}) Tumoro ek osto-tisuo. DE.}
\l{}
\l{\sou{ostito.}\cel{1}(\sou{patol.}) Inflamuro di la osto-tisuo. - DEF.}
\l{}
\l{\sou{ostro.}\cel{2}(\sou{zool.)}\cel{2}Genero de moluski lamelo-brankia, tipo di}
\l{\cel{3}la familio \textquotedbl{}ostreidi\textquotedbl{}, ek cirkume cent speci qui vivas}
\l{\cel{3}en la mari varma e mi-varma, e kinacent speci fosila}
\l{\cel{3}en la strati sekundara e terciara. - L.\cel{1}\sou{ostrea edulis}}
\l{\cel{3}e\cel{1}\sou{gryphea angulata}. - DEFIRS.}
}
